% Korean Support Test Document for XeLaTeX
\documentclass[journal]{IEEEtran}

% Import XeLaTeX configuration for Korean support
% XeLaTeX Configuration for Korean Support
% Include this file in your LaTeX document preamble

% --- XeLaTeX Packages ---
\usepackage{xeCJK}
\usepackage{fontspec}

% --- Font Configuration ---
% Main font family (English)
\setmainfont[Ligatures=TeX]{Times New Roman}
\setsansfont[Ligatures=TeX]{Arial}
\setmonofont{Courier New}

% CJK font configuration (Korean, Chinese, Japanese)
% Using Noto Sans CJK which provides excellent Korean support
\setCJKmainfont[Scale=1]{Noto Sans CJK KR}
\setCJKsansfont[Scale=1]{Noto Sans CJK KR}
\setCJKmonofont[Scale=0.9]{Noto Sans Mono CJK KR}

% --- Alternative Korean Fonts (if Noto Sans CJK is not available) ---
% Uncomment one of the following based on your installed fonts:
% \setCJKmainfont[Scale=1]{SimSun}          % Chinese (alternative)
% \setCJKmainfont[Scale=1]{Liberation Sans} % Generic (fallback)

% --- Additional CJK Settings ---
% Proper spacing and line breaks for CJK text
\XeTeXlinebreaklocale "zh_Hans"
\XeTeXlinebreakskip = 0pt plus 1pt minus 0.5pt


% --- Other Packages ---
\usepackage{cite}
\usepackage{amsmath,amssymb,amsfonts}
\usepackage{algorithmic}
\usepackage{graphicx}
\usepackage{textcomp}
\usepackage{xcolor}
\usepackage{hyperref}

% --- Document Info ---
\title{An Integrated Industrial Inspection System based on Deep Learning and Geometric Constraints}
\subtitle{한국어 지원 테스트 (Korean Support Test)}

\author{Author Name
\thanks{This work was supported by ...}}

\begin{document}

\maketitle

\begin{abstract}
This paper presents an integrated industrial inspection system designed for automated equipment diagnosis.

한국어 지원: 본 논문은 깊은 학습과 기하학적 제약 조건을 기반으로 한 통합 산업 검사 시스템을 제시합니다.
\end{abstract}

\begin{IEEEkeywords}
Industrial Inspection, Deep Learning, Computer Vision, Gauge Reading, Automated Diagnosis, 산업검사, 딥러닝, 컴퓨터비전
\end{IEEEkeywords}

\section{Introduction}
\label{sec:intro}

This section introduces the research problem.

한국 소개: 본 섹션은 연구 문제를 소개합니다.

\section{Methodology}
\label{sec:method}

The proposed methodology combines YOLO detection with geometric calculations.

제안된 방법론은 YOLO 감지와 기하학적 계산을 결합합니다.

\section{Experiments}
\label{sec:exp}

The experimental results demonstrate the effectiveness of the proposed approach.

실험 결과는 제안된 접근 방식의 효과를 보여줍니다.

\section{Conclusion}
\label{sec:conclusion}

This paper presented an integrated solution for industrial inspection systems with Korean language support.

본 논문은 한국어 언어 지원이 있는 산업 검사 시스템을 위한 통합 솔루션을 제시했습니다.

\begin{thebibliography}{1}
\bibitem{key1} Reference example
\end{thebibliography}

\end{document}
